% ============================================================
% TABLE 1 — Main Results (with Who&When baselines)
% Extends paper/main_results_table.tex by adding two Who&When
% localization baselines (W1, W2) on the open-source panel.
%
% Best method selected: +GI+SI (graph injection + span index)
%   — best on Gemini-GAIA and QwenLong-GAIA(dedup-C)
%   — ablation in Table 2 justifies this choice
%
% Dataset conventions:
%   Gemini-Flash : original split (zero_shot/), full traces
%   Gemini-Pro   : original split (zero_shot2/), full traces
%   Mistral / GPT-120B / GPT-20B / Gemma / QwenLong :
%                  GAIA_dedup / SWE_Bench_dedup, full traces  [†]
%
% Who&When rows (W1, W2) live under each open-source model.
%   W1: holistic localization — one call asks the model to label
%       every span in the trace at once.
%   W2: step-by-step localization — one call per span; the prompt
%       accumulates a prefix of prior spans.
% Source: baselines/outputs/outputs_*-{w1,w2}-metrics.txt
%   (run via baselines/who_and_when/run_who_and_when_vllm.py).
% Gemini-Flash/Pro intentionally excluded from Who&When (see
%   baselines/TODO.md §8 — cost).
%
% NOTE (framing): Who&When is a dataset + prompting recipe, not a
% competing method on TRAIL. The Baseline row in the original
% main_results_table.tex is already the single→multi-error
% adaptation of the W&W prompt. W1/W2 rows here are kept as a
% working draft for a future localization-strategy ablation
% (pure W1/W2 vs. W1+graph / W2+graph), not for the headline
% main results table.
%
% Dedup policy: dedup splits used only for small/open models;
%   Gemini (large frontier) uses orig split only.
%
% Scores recalculated with fixed evaluator: errors now counted as empty
% predictions instead of being skipped.
% ============================================================

\begin{table*}[t]
\centering
\small
\setlength{\tabcolsep}{5pt}
\renewcommand{\arraystretch}{1.2}
\begin{tabular}{l l rrr @{\hspace{10pt}} rrr}
\toprule
\multirow{2}{*}{\textbf{Model}} &
\multirow{2}{*}{\textbf{Method}} &
\multicolumn{3}{c}{\textbf{TRAIL-GAIA}} &
\multicolumn{3}{c}{\textbf{TRAIL-SWE-Bench}} \\
\cmidrule(lr){3-5}\cmidrule(lr){6-8}
& & W-F1 & Loc & Joint & W-F1 & Loc & Joint \\
\midrule
\multirow{2}{*}{Gemini-2.5-Flash}
  & Baseline  & 37.07 & 34.22 & 12.71 &  8.71 &  0.63 & 0.00 \\
  & +GI+SI    & \textbf{40.75} & \textbf{35.14} & \textbf{14.48} &
               \textbf{21.56} & \textbf{3.44} & 0.00 \\[3pt]
\multirow{2}{*}{Gemini-2.5-Pro}
  & Baseline  & 33.51 & 35.68 & 13.74 &
               26.66 &  6.42 &  1.35 \\
  & +GI+SI    & \textbf{40.19} & \textbf{41.12} & 13.54 &
               \textbf{38.77} & \textbf{16.75} & \textbf{3.34} \\[3pt]
\multirow{4}{*}{Mistral-Small-3.1-24B$^{\dagger}$}
  & Baseline  & 24.06 & 23.83 &  3.78 &
               9.80 &  9.36 &  1.57 \\
  & W1        & 22.41 & \textbf{27.59} & \textbf{4.20} &
               \textbf{24.81} &  6.78 &  0.58 \\
  & W2        & 14.74 & \textbf{25.06} &  2.99 &
                3.93 &  1.89 &  0.28 \\
  & +GI+SI    & \textbf{30.76} & \textbf{27.79} & \textbf{11.51} &
               6.76 &  1.67 &  0.00 \\[3pt]
\multirow{4}{*}{GPT-oss-120B$^{\dagger}$}
  & Baseline  & 25.76 & 16.92 &  4.59 &
               22.08 &  2.58 &  0.00 \\
  & W1        & 15.87 & \textbf{27.19} &  3.81 &
               19.10 & \textbf{5.35} & \textbf{0.56} \\
  & W2        & 22.72 & \textbf{60.52} &  2.72 &
               13.22 & \textbf{37.91} & \textbf{0.88} \\
  & +GI+SI    & \textbf{30.14} & \textbf{24.47} & \textbf{6.34} &
               \textbf{27.13} & \textbf{2.92} &  0.00 \\[3pt]
\multirow{4}{*}{GPT-oss-20B$^{\dagger}$}
  & Baseline  & 23.28 &  6.29 &  2.61 &
               9.31 &  0.50 &  0.00 \\
  & W1        & 13.88 & \textbf{23.72} & \textbf{3.59} &
                8.32 & \textbf{1.61} &  0.00 \\
  & W2        & 20.34 & \textbf{38.65} & \textbf{3.82} &
               \textbf{33.26} & \textbf{34.94} & \textbf{2.16} \\
  & +GI+SI    & 22.29 & \textbf{18.29} & \textbf{4.40} &
               \textbf{13.03} & \textbf{1.33} & \textbf{0.42} \\[3pt]
\multirow{4}{*}{Gemma-3-27B-IT$^{\dagger}$}
  & Baseline  & 16.25 &  1.79 &  0.33 &
               11.39 &  0.21 &  0.00 \\
  & W1        & 14.75 & \textbf{12.61} & \textbf{2.69} &
               \textbf{12.90} & \textbf{6.41} & \textbf{0.25} \\
  & W2        & \textbf{20.18} & \textbf{22.26} & \textbf{0.97} &
               10.49 & \textbf{11.72} & \textbf{0.31} \\
  & +GI+SI    & \textbf{25.30} & \textbf{11.00} & \textbf{1.33} &
               \textbf{15.16} & \textbf{2.36} &  0.00 \\[3pt]
\multirow{4}{*}{QwenLong-L1-32B$^{\dagger}$}
  & Baseline  & 12.40 &  3.92 &  0.27 &
               2.14 &  0.00 &  0.00 \\
  & W1        & 11.07 & \textbf{16.10} & \textbf{1.03} &
               \textbf{5.75} & \textbf{3.80} & \textbf{0.28} \\
  & W2        & \textbf{17.10} & \textbf{18.42} & \textbf{2.98} &
                0.00 & \textbf{31.63} &  0.00 \\
  & +GI+SI    & \textbf{16.83} & \textbf{19.23} & \textbf{3.63} &
               \textbf{11.09} & \textbf{0.42} & 0.00 \\
\bottomrule
\end{tabular}
\caption{%
  Main results on TRAIL-GAIA and TRAIL-SWE-Bench, extended with
  Who\&When localization baselines.
  \textbf{W-F1}: weighted category F1.
  \textbf{Loc}: average location accuracy.
  \textbf{Joint}: location-category joint accuracy (all in \%).
  \textbf{W1}: Who\&When holistic localization (one call labels all
  spans together).
  \textbf{W2}: Who\&When step-by-step localization (one call per span
  with cumulative-prefix prompts).
  \textbf{+GI+SI}: causal graph injection with span-index prefixes
  (our method).
  Bold: improvement over the model's own TRAIL Baseline row.
  $^{\dagger}$Uses the deduplicated split (\texttt{GAIA\_dedup} /
  \texttt{SWE\_Bench\_dedup}).
  Gemini-Flash/Pro are excluded from Who\&When (W1/W2) due to
  per-trace call-count cost (see appendix).
  \text{---}: experiment not run.%
}
\label{tab:main_results_baseline}
\end{table*}


% ============================================================
% TABLE 2 — Ablation: Method Variants
% Gemini-2.5-Flash on TRAIL-GAIA (original split).
% Justifies the choice of +GI+SI as the primary method.
% N: Baseline/+CG = 109; +CG+SI = 106; +GI+SI = 110.
% ============================================================

\begin{table}[t]
\centering
\small
\setlength{\tabcolsep}{6pt}
\renewcommand{\arraystretch}{1.2}
\begin{tabular}{l rrr}
\toprule
\textbf{Method} & \textbf{W-F1} & \textbf{Loc} & \textbf{Joint} \\
\midrule
Baseline         & 37.07 & 34.22 & 12.71 \\
+CG              & 40.21 & \textbf{35.14} & 14.53 \\
+CG+SI           & 39.60 & 34.40 & \textbf{15.81} \\
+GI+SI           & \textbf{40.75} & \textbf{35.14} & 14.48 \\
\bottomrule
\end{tabular}
\caption{%
  Ablation of causal-graph variants on TRAIL-GAIA
  (Gemini-2.5-Flash, original split, all metrics in \%).
  \textbf{+CG}: causal graph edges in prompt.
  \textbf{+CG+SI}: causal graph with span-index prefixes.
  \textbf{+GI+SI}: graph injection with span-index prefixes (our method).
  Bold: best per metric.%
}
\label{tab:ablation_methods_baseline}
\end{table}


% ============================================================
% TABLE 3 — Ablation: Deduplicated vs Original Split
% Mistral-Small-3.1-24B on TRAIL-GAIA.
% Shows impact of dataset deduplication on evaluation quality.
%   orig  : zero_shot/, original split,  N ≈ 48–49
%   dedup : zero_shot2/, GAIA_dedup,     N ≈ 64–81
% No baseline exists for the original split.
% Dedup baseline exists (zero_shot2/, N=81).
% ============================================================

\begin{table}[t]
\centering
\small
\setlength{\tabcolsep}{6pt}
\renewcommand{\arraystretch}{1.2}
\begin{tabular}{l l rrr}
\toprule
\textbf{Split} & \textbf{Method} & \textbf{W-F1} & \textbf{Loc} & \textbf{Joint} \\
\midrule
\multirow{2}{*}{Original}
  & +CG      & 17.98 &  8.69 &  2.51 \\
  & +CG+SI   & 20.85 & 14.59 &  3.09 \\
\midrule
\multirow{4}{*}{Dedup}
  & Baseline & 24.06 & 23.83 &  3.78 \\
  & +CG      & 28.31 & \textbf{30.31} & 11.16 \\
  & +CG+SI   & 29.25 & 25.99 &  8.75 \\
  & +GI+SI   & \textbf{30.76} & 27.79 & \textbf{11.51} \\
\bottomrule
\end{tabular}
\caption{%
  Effect of dataset deduplication on TRAIL-GAIA evaluation
  (Mistral-Small-3.1-24B, all metrics in \%).
  \textit{Original}: zero-shot run on the full (non-deduplicated) split ($N{\approx}48$).
  \textit{Dedup}: run on the deduplicated \texttt{GAIA\_dedup} split ($N{\approx}73$).
  No GAIA baseline exists for the original split.
  Bold: best per metric across all rows.%
}
\label{tab:ablation_dedup_baseline}
\end{table}


% ============================================================
% TABLE 4 — Trace Coverage: N per model × method × dataset
% Shows how many traces were actually processed in each run.
% N comes from the correlation-table "N" column in each metrics file.
% Who&When (W1/W2) N values from baselines/outputs/*-metrics.txt.
% Key observations:
%   - SWE N is small (9–31) for all models → wide CI on SWE numbers
%   - Mistral +GI+SI GAIA N=64 vs. Baseline N=81 → incomplete run
%   - GPT-20B Baseline GAIA N=73 vs. +CG N=86 → batched rollout
%   - QwenLong SWE +CG+SI N=16 (lowest) → context overflow
%   - QwenLong W2 GAIA N=37 → context overflow on long step-by-step prompts
% ============================================================

\begin{table}[t]
\centering
\small
\setlength{\tabcolsep}{5pt}
\renewcommand{\arraystretch}{1.1}
\begin{tabular}{ll rrrrrr}
\toprule
\textbf{Model} & \textbf{Dataset} & \textbf{Base} & \textbf{W1} & \textbf{W2} & \textbf{+CG} & \textbf{+CG+SI} & \textbf{+GI+SI} \\
\midrule
Gemini-2.5-Flash      & GAIA (orig)    & 109 & \text{---} & \text{---} & 109 & 106 & 110 \\
                      & SWE  (orig)    &   9 & \text{---} & \text{---} &  22 &  20 &  20 \\
\midrule
Gemini-2.5-Pro        & GAIA (orig)    & 117 & \text{---} & \text{---} & \text{---} & \text{---} & 110 \\
                      & SWE  (orig)    &  31 & \text{---} & \text{---} & \text{---} & \text{---} &  31 \\
\midrule
Mistral-Small-24B     & GAIA (dedup)   &  81 & 112 & 114 &  76 &  73 &  64 \\
                      & SWE  (dedup)   &  11 &  26 &  29 & \text{---} & \text{---} &   9 \\
\midrule
GPT-oss-120B          & GAIA (dedup)   &  90 & 114 & 114 &  90 &  89 &  90 \\
                      & SWE  (dedup)   &  19 &  31 &  30 &  19 &  19 &  19 \\
\midrule
GPT-oss-20B           & GAIA (dedup)   &  73 &  88 &  89 &  86 &  89 &  87 \\
                      & SWE  (dedup)   &  12 &  30 &  29 &  18 &  19 &  19 \\
\midrule
Gemma-3-27B-IT        & GAIA (dedup)   &  88 & 116 & 116 &  88 &  89 &  87 \\
                      & SWE  (dedup)   &  19 &  30 &  30 &  19 &  19 &  19 \\
\midrule
QwenLong-L1-32B       & GAIA (dedup)   &  88 & 113 &  37 &  87 &  87 &  84 \\
                      & SWE  (dedup)   &  19 &  29 &  30 &  17 &  16 &  18 \\
\bottomrule
\end{tabular}
\caption{%
  Number of traces processed per model, dataset, and method.
  Variation across methods within the same model--dataset pair indicates
  incomplete runs (most pronounced for Mistral +GI+SI: $N{=}64$ vs.\
  baseline $N{=}81$) or context-overflow failures counted as empty
  predictions.
  W1 / W2 trace counts come from the Who\&When runs in
  \texttt{baselines/outputs/}; the higher GAIA $N$ ($\sim$112--116)
  reflects the unfiltered who\_and\_when input set, while TRAIL-prompt
  rows use the agent-trace dedup ($N{=}73$--$90$).
  QwenLong W2 GAIA reaches only $N{=}37$ due to context-overflow failures
  on the cumulative-prefix prompts that grow $O(N_{\text{spans}}^2)$.
  QwenLong GAIA dedup runs are sourced from \texttt{zero\_shot/compressed/GAIA\_dedup-*}
  (full-trace dedup experiments); the \texttt{GAIA\_compressed-*} runs
  in the same folder use a different context-compression method and are omitted.
  SWE $N$ is small across all models ($9$--$31$), making SWE comparisons
  indicative rather than definitive.
  GPT-oss-20B GAIA baseline ($N{=}73$) is lower than its graph conditions
  ($N{=}86$--$89$) due to batched rollout across dataset versions.
  QwenLong SWE +CG+SI reaches the minimum $N{=}16$, consistent with the
  span-index token overhead causing more context-overflow failures on the
  longer SWE traces.
  \text{---}: condition not run.%
}
\label{tab:trace_coverage_baseline}
\end{table}
